\rhead{EE101: EE Fundamentals}
\phantomsection 
\section*{EE Fundamentals (EE101)}
\label{course:EE101}

\noindent \small{\hyperref[main:scheme]{$\leftarrow$ Back to Scheme}} \vspace{0.5cm}

\noindent \textbf{Credits:} 3 (3-0-0)

\subsection*{Course Content}
\noindent\textbf{Unit 1} \\
\textit{DC Circuit Analysis and Power Foundations:} Fundamental Concepts: Electrical Current, Voltage, Resistance, Power, and Energy. The Concept of "Ground" and Reference Voltages in Digital Systems. Ohm's Law. Circuit Elements: Resistors, Inductors, and Capacitors. Series and Parallel Combinations. Practical concepts of Pull-up and Pull-down Resistors. Analysis Methods: Kirchhoff's Laws (KCL/KVL), Mesh Analysis, and Nodal Analysis. Network Theorems: Thevenin’s, Norton’s, and Maximum Power Transfer Theorem.

\vspace{0.3cm}

\noindent\textbf{Unit 2} \\
\textit{AC Circuits, Transients, and Signal Basics:} AC Fundamentals: Generation of Sinusoidal Voltage, Frequency, Period, Phase, Average, and RMS Values. Introduction to Non-Sinusoidal Signals (Square Waves and Pulses). Transient Analysis: Behavior of RC and RL circuits with DC excitation, Time Constants, and their relation to Digital Signal Propagation Delay. AC Circuit Analysis: Phasor Representation, Series R-L, R-C, and R-L-C Circuits. Impedance, Resonance, and Bandwidth. Signal Conditioning: Passive Low-Pass and High-Pass RC Filters.

\vspace{0.3cm}

\noindent\textbf{Unit 3} \\
\textit{Semiconductor Devices and Digital Physics:} Diodes: P-N Junction Physics, V-I Characteristics. Applications: Rectifiers, Zener Regulation, and Clipping/Clamping Circuits. Bipolar Junction Transistor (BJT): Operation as a Switch (Cutoff vs. Saturation). Field-Effect Transistors (MOSFET): JFET and Enhancement-mode MOSFETs. Comparison of BJT vs. MOSFET logic. CMOS Technology: The CMOS Inverter and physical construction of Logic Gates (AND, OR, NOT).

\vspace{0.3cm}

\noindent\textbf{Unit 4} \\
\textit{Operational Amplifiers and Data Conversion:} Operational Amplifiers (Op-Amps): Characteristics of Ideal vs. Practical Op-Amps. Virtual Ground Concept. Configurations: Inverting, Non-Inverting, and Unity Gain Buffer. Applications: Summing Amplifier, Differentiator, Integrator, and Comparator (Zero Crossing Detector). Data Conversion: Analog-to-Digital Converters (ADC) and Digital-to-Analog Converters (DAC). Sampling Theorem (Nyquist Rate) and Quantization.

\vspace{0.3cm}

\noindent\textbf{Unit 5} \\
\textit{Electronic Systems, Interfacing, and Communication:} Sensors and Transducers: Active vs. Passive Sensors (LDR, Photodiode, Thermistor, Ultrasonic). Signal Conditioning requirements. Actuators and Drivers: Interfacing Microcontrollers to High-Power Loads (Relays, DC Motors). Opto-isolators for Galvanic Isolation. H-Bridge Circuits for Motor Direction Control. PWM (Pulse Width Modulation) basics.

\newpage
\rhead{Curriculum Schema}